
\subsubsection{5.1 RQ1: Gap Learnability and Target Independence}

Figure 1 shows Spearman rank correlations for all four encoders on the gap prediction target. Two encoders
exceed the existence threshold (r > 0.5): FlatMLP (r = 0.5567 in the existence run, h-e1) and DWSNet
(r = 0.5104 in h-e1). This establishes that generalization gap is a learnable signal in weight tensors
--- accessible even to a position-indexed flat encoder without equivariant constraints.

\textbf{Table 1: Gap and Test Accuracy Spearman (Test Split, N = 1,000)}

\begin{table*}[htbp]
\centering
\resizebox{\dimexpr 2\columnwidth+\columnsep\relax}{!}{%
\begin{tabular}{lllll}
\toprule
Encoder & Spearman(gap) & 95\% CI & Spearman(test\_acc) & 95\% CI \\
\midrule
FlatMLP & 0.5567† & [0.4850, 0.5801]‡ & 0.2790 & [0.2173, 0.3343] \\
DWSNet & 0.5104† & [0.4377, 0.5325]‡ & 0.4553 & [0.4020, 0.5054] \\
NFT & 0.5752 & [0.5339, 0.6158] & 0.4801 & [0.4326, 0.5262] \\
GNN & 0.3747 & [0.3180, 0.4265] & 0.3480 & [0.2913, 0.4037] \\
\bottomrule
\end{tabular}
}
\end{table*}

*† FlatMLP and DWSNet gap point estimates in Table 1 are from the existence run (h-e1). FlatMLP gap in
the architecture comparison run (h-m1) is 0.5330; this h-m1 value is used as the control baseline in the
$\Delta$ computation (Table 2), consistent with the run that produced the equivariant encoder CIs. Run-to-run
gap variance for FlatMLP (0.5567 vs. 0.5330, ~4.4\%) reflects the 3-trial search budget.*

*‡ FlatMLP CI [0.4850, 0.5801] and DWSNet CI [0.4377, 0.5325] are from h-m1 bootstrap (N = 1,000
resamples); shown here alongside NFT CI [0.5339, 0.6158] for comparison.*

The A1 audit confirms that gap is not trivially derivable from test accuracy: Spearman(gap, $-test_acc)$
= $-0.142$, far below the collinearity threshold of 0.95. Predicting gap is a genuinely independent task.

\textbf{The counterintuitive asymmetry.} Figure 2 reveals that FlatMLP predicts gap (r = 0.557) substantially
better than test accuracy (r = 0.279) from identical weight inputs. This reverses the expectation from
prior literature (where FlatMLP achieves r $\approx$ 0.85 on test\_acc in the original Unterthiner zoo). We
discuss this anomaly in Section 6; it does not affect the gap prediction results, which are our
primary contribution.

\subsubsection{5.2 RQ2: Architecture Ranking on Gap Prediction}

\textbf{NFT achieves the highest gap Spearman} (r = 0.5752) with a 95\% CI of [0.5339, 0.6158]. FlatMLP's
95\% CI from the same run (h-m1) is [0.4850, 0.5801]. The two intervals overlap, but NFT's lower bound
(0.5339) exceeds FlatMLP's point estimate (0.5330), providing directional evidence of NFT's advantage;
the CI overlap prevents a non-overlap claim. The ranking is: NFT > FlatMLP > DWSNet > GNN for gap
prediction.

The architecture-specificity of this result is a key finding. DWSNet (within-layer equivariance)
achieves r = 0.4881 on gap --- below FlatMLP's r = 0.5330. GNN achieves only r = 0.3747, the lowest
among all encoders. Equivariance is not uniformly helpful for gap prediction: only the cross-layer
attention architecture (NFT) provides a meaningful advantage, while within-layer and graph-structured
equivariance underperform the non-equivariant baseline.

Figure 2 (predicted vs. true gap scatter, h-m1) shows the qualitative fit across encoders on the test
split. NFT's scatter is tighter and more linear than DWSNet's or GNN's, consistent with its higher
Spearman correlation.

\subsubsection{5.3 RQ3: Gap-Specific Signal Independent of Test Accuracy (P3)}

The partial Spearman analysis tests whether NFT's gap predictions contain information about true gap
that cannot be explained by test accuracy rank alone.

\textbf{NFT partial Spearman (gap | test\_acc): r = 0.7305, p = 1.$60\times 10^{-167}$}

Figure 3 shows the residual scatter for the P3 analysis. After partialling out test accuracy rank
from both NFT's gap predictions and the true gap labels, the residual correlation is r = 0.73.
This result has two noteworthy properties:

\begin{enumerate}
\item \textbf{Effect size relative to direct Spearman.} The partial correlation (r = 0.73) is larger in
\end{enumerate}
   magnitude than the direct Spearman (r = 0.57). Noting that partial and direct Spearman measure
   different quantities and are not directly statistically comparable, this pattern suggests that
   the gap-specific component of NFT's predictions is particularly strong relative to the combined
   signal, and that partialling out test accuracy rank isolates an overfitting-related structure
   that NFT captures well.

\begin{enumerate}
\item \textbf{Unambiguous significance.} With N = 1,000 test samples, p = 1.$60\times 10^{-167}$ is not a borderline
\end{enumerate}
   result --- the gap-specific signal is substantial and reproducible.

This result answers RQ3: NFT gap predictions contain information about true generalization gap that is
statistically independent of test accuracy. Generalization gap encodes a distinct signal in weight space
that cross-layer attention is well-positioned to extract.

\subsubsection{5.4 RQ4: Differential Advantage ($\Delta )$ --- Null Result}

Figure 4 and the bootstrap CI figure show $\Delta$ values for equivariant encoders with 95\% bootstrap CIs.

\textbf{Table 2: Differential Advantage $\Delta$ = [gap improvement] $-$ [test\_acc improvement] over FlatMLP}

\begin{table*}[htbp]
\centering
\resizebox{\dimexpr 2\columnwidth+\columnsep\relax}{!}{%
\begin{tabular}{llllll}
\toprule
Encoder & gap Spearman & acc Spearman & $\Delta$ & 95\% CI & $\Delta$ > 0.02? \\
\midrule
DWSNet & 0.4881 & 0.4553 & $-0.2212$ & [$-0.298, -0.150]$ & No \\
NFT & 0.5752 & 0.4801 & $-0.1589$ & [$-0.231, -0.087]$ & No \\
GNN & 0.3747 & 0.3480 & $-0.2272$ & [$-0.323, -0.131]$ & No \\
\bottomrule
\end{tabular}
}
\end{table*}

\textit{Gate criterion: $\Delta$ > 0.02 for $\geq 2$ encoders. Result: 0/3 encoders pass. Gate: FAIL.}
\textit{FlatMLP baselines used: gap = 0.5330 (h-m1), test\_acc = 0.2790 (h-m2), both from the same experimental run.}

All equivariant $\Delta$ values are strongly negative, with 95\% CIs entirely below zero. The differential
advantage hypothesis is not confirmed under these experimental conditions.

\textbf{Why are $\Delta$ values negative?} The dominant driver is the test\_acc improvement component --- all
equivariant encoders show substantially larger test\_acc improvement over FlatMLP than gap improvement.
Since FlatMLP test\_acc (r = 0.279) is anomalously low (vs. literature ~0.85), the apparent acc
improvement for equivariant encoders is inflated. DWSNet shows 0.455 $-$ 0.279 = +0.176 improvement
on test\_acc, while its gap improvement is 0.488 $-$ 0.533 = $-0.045$ (DWSNet is actually below FlatMLP
on gap). This confound prevents a clean interpretation of the $\Delta$ result.

\textbf{What we can conclude:} (1) Under our experimental conditions, equivariant encoders do not show a
gap-specific differential advantage. (2) The most likely confounder is the FlatMLP test\_acc baseline
being severely underestimated due to our limited search budget and potential zoo generation differences.
(3) The gap prediction results (RQ1-RQ3) are not affected by this issue --- they rely only on gap
Spearman, which is internally consistent across independent runs (FlatMLP: 0.5567 in h-e1, 0.5330
in h-m1; deviation < 5\%).

